\documentclass{csexam}

\usepackage{listings}
\usepackage{xcolor}
\usepackage{tikz}

% Python code styling
\lstset{
    language=Python,
    basicstyle=\ttfamily\small,
    keywordstyle=\color{blue}\bfseries,
    stringstyle=\color{red},
    commentstyle=\color{green!50!black}\itshape,
    showstringspaces=false,
    breaklines=true,
    frame=single,
    numbers=left,
    numberstyle=\tiny\color{gray}
}

% Course information
\university{St. Francis Xavier University}
\department{Department of Computer Science}
\course{CSCI 128: Coding for Problem Solving}
\courseshort{CSCI 128}
\exam{Midterm 1 - SOLUTIONS}
\examdate{Winter 2026}

\begin{document}

\maketitle
\printanswers

\begin{questions}

% ======================================
% Section 1: Values, Types, and Variables
% ======================================
\question[6] For each expression below, determine the data type of the result. If there is an error, write "ERROR".

\begin{solution}
\begin{parts}
    \part \texttt{5 + 3.0} \hspace{0.5cm} Type: \textbf{float}
    \part \texttt{str(42) + "7"} \hspace{0.5cm} Type: \textbf{str}
    \part \texttt{int("3.14")} \hspace{0.5cm} Type/Result: \textbf{ERROR} (ValueError - cannot convert string with decimal to int)
    \part \texttt{10 / 2} \hspace{0.5cm} Type: \textbf{float}
    \part \texttt{10 // 3} \hspace{0.5cm} Type: \textbf{int}
    \part \texttt{"Python" * 2} \hspace{0.5cm} Type: \textbf{str}
\end{parts}
\end{solution}

\question[8] What is the output of the following code? Show exactly what gets printed.

\begin{lstlisting}
x = "5"
y = 3
x = int(x) + y
y = str(y) * 2
print(x, y)
\end{lstlisting}

\begin{solution}
\textbf{Output:}
\begin{verbatim}
8 33
\end{verbatim}

\textbf{Explanation:}
\begin{itemize}
    \item Line 1: x = "5" (string)
    \item Line 2: y = 3 (int)
    \item Line 3: x = int("5") + 3 = 5 + 3 = 8
    \item Line 4: y = str(3) * 2 = "3" * 2 = "33"
    \item Line 5: prints 8 and "33" separated by a space
\end{itemize}
\end{solution}

\question[8] Trace through the following code carefully and show what gets printed:
\begin{lstlisting}
a = 10
b = 20
a = b - a
b = a + b
a = b - a
print(a)
print(b)
\end{lstlisting}

\begin{solution}
\textbf{Output:}
\begin{verbatim}
20
10
\end{verbatim}

\textbf{Trace:}
\begin{itemize}
    \item Line 1: a = 10
    \item Line 2: b = 20
    \item Line 3: a = 20 - 10 = 10
    \item Line 4: b = 10 + 20 = 30
    \item Line 5: a = 30 - 10 = 20
    \item Line 6: prints 20
    \item Line 7: prints 10
\end{itemize}

Note: This is a classic swap algorithm using arithmetic operations.
\end{solution}

\newpage

\question[14] Write a complete Python program that:
\begin{itemize}
    \item Asks the user for their age
    \item Asks the user for the number of years to add
    \item Calculates their future age
    \item Calculates how many months old they will be (future age × 12)
    \item Prints both results with appropriate labels
\end{itemize}

\begin{solution}
\begin{lstlisting}
# Get user input
age = int(input("Enter your age: "))
years = int(input("Enter number of years: "))

# Calculate future age and months
future_age = age + years
months = future_age * 12

# Print results
print("In", years, "years you will be:", future_age)
print("That is", months, "months old")
\end{lstlisting}

\textbf{Alternative solution using f-strings:}
\begin{lstlisting}
# Get user input
age = int(input("Enter your age: "))
years = int(input("Enter number of years: "))

# Calculate future age and months
future_age = age + years
months = future_age * 12

# Print results
print(f"In {years} years you will be: {future_age}")
print(f"That is {months} months old")
\end{lstlisting}
\end{solution}

\newpage

\question[8] What is wrong with each of the following code snippets? Identify the error and explain how to fix it.

\begin{solution}
\begin{parts}
    \part
\begin{lstlisting}
age = input("Enter your age: ")
next_year = age + 1
print(next_year)
\end{lstlisting}
    \textbf{Error:} \textit{Type error - cannot add string and integer. input() returns a string.}

    \textbf{Fix:} \textit{Convert input to int: age = int(input("Enter your age: "))}

    \part
\begin{lstlisting}
name = input("Enter your name: ")
print("Hello " + name + ". You are " + 25 + " years old")
\end{lstlisting}
    \textbf{Error:} \textit{Type error - cannot concatenate string and integer (25).}

    \textbf{Fix:} \textit{Convert 25 to string: print("Hello " + name + ". You are " + str(25) + " years old")}
\end{parts}
\end{solution}

\question[6] Explain the difference between the operators \texttt{/} and \texttt{//} in Python. Give an example showing different results for each.

\begin{solution}
\textbf{Difference:}
\begin{itemize}
    \item \texttt{/} is \textbf{true division} - always returns a float with the exact result including decimal places
    \item \texttt{//} is \textbf{floor division} (or integer division) - returns an integer by discarding the decimal portion (rounds down)
\end{itemize}

\textbf{Examples:}
\begin{itemize}
    \item \texttt{10 / 3} results in \texttt{3.3333...} (float)
    \item \texttt{10 // 3} results in \texttt{3} (int)
    \item \texttt{10 / 2} results in \texttt{5.0} (float, even though evenly divisible)
    \item \texttt{10 // 2} results in \texttt{5} (int)
\end{itemize}
\end{solution}

\newpage

% ======================================
% Section 2: Pictures and Pixels
% ======================================
\question[6] Answer the following questions:

\begin{solution}
\begin{parts}
    \part \textbf{(128, 128, 128):} Medium gray. \textbf{(200, 200, 200):} Light gray (closer to white).

    When all three RGB values are equal, you get a shade of gray. Lower values are darker, higher values are lighter.

    \part \textbf{Yellow:} (255, 255, 0)

    Yellow is created by mixing red and green at full intensity with no blue.

    \part \textbf{Orange:} Approximately (255, 165, 0) or similar values like (255, 128, 0)

    Orange is red at full intensity with a moderate amount of green and no blue.
\end{parts}
\end{solution}

\question[5] For an image with dimensions 100×100:

\begin{solution}
\begin{parts}
    \part \textbf{Top-right corner:} (99, 0)

    x ranges from 0 to 99, y ranges from 0 to 99. Top means y=0, right means x=99.

    \part \textbf{Bottom-left corner:} (0, 99)

    Bottom means y=99, left means x=0.

    \part \textbf{Center pixel:} (50, 50) or (49, 49)

    Either answer is acceptable since 100×100 doesn't have a true "center" pixel. (50, 50) is slightly off-center but commonly used.
\end{parts}
\end{solution}

\newpage

\question[15] Write a complete Python program using Pillow that:
\begin{itemize}
    \item Loads an image file named \texttt{"photo.jpg"}
    \item Gets the width and height of the image and prints them
    \item Gets the pixel color at position (10, 10)
    \item Extracts the individual R, G, and B values from that pixel
    \item Calculates the sum of the three color values (R + G + B)
    \item Prints the results in the specified format
\end{itemize}

\begin{solution}
\begin{lstlisting}
from PIL import Image

# Load the image
img = Image.open("photo.jpg")

# Get width and height
width, height = img.size

# Print image size
print("Image size:", width, "x", height)

# Get pixel color at (10, 10)
pixel_color = img.getpixel((10, 10))

# Extract R, G, B values
r, g, b = pixel_color

# Calculate sum
total = r + g + b

# Print results
print("Pixel at (10, 10) - Red:", r, ", Green:", g, ", Blue:", b)
print("Total color value:", total)
\end{lstlisting}

\textbf{Alternative with f-strings:}
\begin{lstlisting}
from PIL import Image

img = Image.open("photo.jpg")
width, height = img.size

print(f"Image size: {width} x {height}")

pixel_color = img.getpixel((10, 10))
r, g, b = pixel_color
total = r + g + b

print(f"Pixel at (10, 10) - Red: {r}, Green: {g}, Blue: {b}")
print(f"Total color value: {total}")
\end{lstlisting}
\end{solution}

\newpage

\question[12] Complete the following Python program by filling in the blanks. This program creates a vertical red stripe down the middle of a 200×200 image:

\begin{solution}
\begin{lstlisting}
from PIL import Image

# Create a new white image (200x200)
img = Image.new("RGB", (200, 200), (255, 255, 255))

# Draw a red vertical stripe in the middle (x from 95 to 105)
for x in range(95, 105):
    for y in range(200):
        img.putpixel((x, y), (255, 0, 0))

# Save the image
img.save("stripe.jpg")
\end{lstlisting}

\textbf{Explanation of blanks:}
\begin{itemize}
    \item \texttt{Image.new} - method to create a new image
    \item \texttt{200, 200} - dimensions (width, height)
    \item \texttt{95, 105} - range for x coordinate (middle stripe)
    \item \texttt{200} - range for y coordinate (full height, 0 to 199)
    \item \texttt{img.putpixel} - method to set pixel color
    \item \texttt{255, 0, 0} - red color (R=255, G=0, B=0)
    \item \texttt{img.save} - method to save the image
\end{itemize}
\end{solution}

\question[12] Write a complete Python program using Pillow that:
\begin{itemize}
    \item Creates a new 100×100 RGB image with all pixels initially black
    \item Uses nested loops to set all pixels in the left half (x from 0 to 49) to red (255, 0, 0)
    \item Uses nested loops to set all pixels in the right half (x from 50 to 99) to green (0, 255, 0)
    \item Saves the result as \texttt{"two\_colors.jpg"}
\end{itemize}

\begin{solution}
\begin{lstlisting}
from PIL import Image

# Create a new 100x100 black image
img = Image.new("RGB", (100, 100), (0, 0, 0))

# Set left half to red
for x in range(0, 50):
    for y in range(100):
        img.putpixel((x, y), (255, 0, 0))

# Set right half to green
for x in range(50, 100):
    for y in range(100):
        img.putpixel((x, y), (0, 255, 0))

# Save the image
img.save("two_colors.jpg")
\end{lstlisting}

\textbf{Alternative solution (more efficient using loops):}
\begin{lstlisting}
from PIL import Image

# Create a new 100x100 black image
img = Image.new("RGB", (100, 100), (0, 0, 0))

# Loop through all pixels
for x in range(100):
    for y in range(100):
        if x < 50:
            img.putpixel((x, y), (255, 0, 0))  # Red
        else:
            img.putpixel((x, y), (0, 255, 0))  # Green

# Save the image
img.save("two_colors.jpg")
\end{lstlisting}
\end{solution}

\end{questions}

\end{document}
